% Short running form; the full title is too wide for the head's right box.
\runninghead{The War God Permits and Bounds}
\properlane{armour-under-providence}{The War God Permits and Bounds: The Adversary on a Leash and the Armour That Is Christ}
\label{sec:armour-under-providence}

Heard from the Introit and the Offertory, the Mass sets two sentences side by
side that seem to contradict each other. The first says that all things are
placed in God's will and that no one can resist it. The second says that Satan
asked to tempt a just man, and that power over his goods and his flesh was given
him by the Lord. Between them the Epistle names the Christian's enemy: not flesh
and blood, but principalities, powers and spirits of wickedness, against whom he
must stand in the evil day clothed in the armour of God. At the end the Communion
asks when God will judge those who persecute the singer. St Gregory the Great
and St John Chrysostom answer the question these texts raise together. Gregory,
expounding Job, finds two wills at work on the one man: the devil's, which
asks for his ruin, and God's, which permits the trial for his proving and keeps
the tempter bridled. Chrysostom, expounding the Epistle, finds the real enemy
behind the human one: the spirits that work in persecutors, against whom alone
the Christian may turn his anger. Read with them, the formulary describes a war
that is fierce and real but bounded. The adversary asks before he strikes and
strikes only as far as he is let; the human persecutor and the offending brother
are not the enemy; the armour is God's, and St Jerome says it is Christ; and the
will that no one resists is the ground of the soldier's confidence, not the
cause of his trouble.

\subsection{Two sentences that seem to collide}

The Introit's sentence is absolute: \latin{In voluntate tua, Domine, universa
sunt posita, et non est qui possit resistere voluntati tuae}. In the Book of
Esther it is said by a man under sentence of death, whose people a royal decree
has handed over to its enemy; he confesses the will that no one resists while
the enemy's will seems to have prevailed. The Offertory tells a story that
sharpens the question. Job is \latin{simplex et rectus, ac timens Deum}, and yet
\latin{Satan petiit, ut tentaret: et data est ei potestas a Domino}. The verb
\latin{petiit} gives the adversary a will of his own, and the passive \latin{data
est \dots\ a Domino} says where his power came from. If no one resists God's will,
how does Satan come to ask for what God's servant does not deserve, and why is
it given? Either the Introit says too much, and the adversary is a second power
whose will God must reckon with; or the Offertory makes God the author of Job's
ruin. The Epistle adds a third term. The Christian has an enemy, and is told not
to mistake him: \latin{non est nobis colluctatio adversus carnem et sanguinem}.
And the Communion, sung after a Gospel that demands forgiveness of the brother,
asks, \latin{quando facies de persequentibus me iudicium?}

\subsection{The petition and the permission: St Gregory the Great}

St Gregory's \work{Moralia} on Job begins with the question why so many strokes
fell on a man who kept every virtue, and answers it with two wills: ``Him the
adversary sought indeed to obtain that he might prove deficient in godliness,
but obtained that he might prove a proficient therein. The Lord in
loving-kindness permitted that to be done, which the devil in his wickedness
required; for when the enemy had got leave to have him with the purpose of
destroying him, by his temptations he effected that his merits were augmented''
(Preface~7). The play in his Latin is on one verb with two prefixes, \latin{ut
deficeret} and \latin{ut proficeret}: the devil sought that Job might fall short,
and obtained that he might go forward. The petition is the devil's and the
permission is God's, and the outcome follows the permission. Job therefore
``became the intermediate subject of the contest between God and the devil''
(Preface~8): the enemy, in putting forth his strength on Job, ``entered the
lists against God''.

On the first permission, ``all that he hath is in thy hand: only put not forth
thy hand upon his person'' (Job 1:12), Gregory bids the reader ``mark in the
Lord's words the dispensations of heavenly pity, how He lets go our enemy, and
keeps him in; how He looses, and yet bridles him. He allows him some things for
temptation, but withholds him from others.'' The Lord gives up Job's substance
and keeps his person, and later gives that too, ``yet He does not loose the enemy
to every thing at once, lest he should crush His own subject by striking him on
every side''. Trials are dealt out by seasons, ``that what by coming all together
would destroy, may when divided be borne up against'', as St Paul says, ``God is
faithful, Who will not suffer you to be tempted above that ye are able'' (1~Cor
10:13). Gregory adds another sense of the same words: the Lord knew ``that His
soldier was brave'', and chose to divide his contests with the enemy so that he
might come out of each ``the more incomparable conqueror'' (\work{Moralia}
II.19). The adversary of the Offertory is on a leash, and the hand that holds it
measures every stroke.

Gregory answers the Introit's claim as well, though from another verse of Job,
``Who hath resisted Him, and had peace?'' (Job 9:4). Resistance to God happens,
he says, but it never succeeds and never has peace: ``whereinsoever there is
resistance made to the Creator, that agreement in peace is broken up''. The
apostate angel, ``banished, has to bear the burthen of himself''; Adam, refusing
obedience, lost even the peace of his body. ``Not that our frailty does resist
His unchangeable decree, but what it has not the power to accomplish, it yet
attempts'' (\work{Moralia} IX.5). Rupert of Deutz, citing the Introit's words
in his treatise on God's omnipotence (he gives them as St Paul's), completes them
from this verse of Job and from Gregory's exposition of it: no one resists God's
will and has peace (\work{De omnipotentia Dei} 9). Read with Gregory, the
Introit does not deny that creatures set themselves against God. It says that
none succeeds, and that the one who tries breaks only his own peace.

\subsection{The will to harm and the power to harm: St Thomas Aquinas}

St Thomas states in a single distinction what Gregory tells in a story. In his
literal exposition of Job he holds that Satan can harm \latin{non quantum vult
sed quantum permittitur}, not as much as he wills but as much as he is permitted:
the will to harm is in every evil being from itself, the power only from God.
When Satan ``went forth from the presence of the Lord'' (Job 1:12), Thomas reads
it as going out from the intention of Providence while unable to escape its
order (\work{Expositio super Iob} 1). In the \work{Summa} he takes the Introit's
own verse as his authority. Asking whether anything can happen outside the order
of the divine government, he answers from Esther: ``It is written (Esther 13:9):
`O Lord, Lord, almighty King, all things are in Thy power, and there is none that
can resist Thy will'''. An effect can fall outside the order of some particular
cause, but not outside the order of the universal cause, and ``as God is the
first universal cause, not of one genus only, but of all being in general, it is
impossible for anything to occur outside the order of the Divine government''
(I q.~103 a.~7).

The two Doctors make the two sentences of the Mass agree without softening
either. The Offertory's \latin{petiit} is the adversary's will, which is his own
and evil; its \latin{data est \dots\ a Domino} is his power, which is God's and
measured. The Introit's will that no one resists is the order to which even the
petition belongs. Gregory speaks of lost peace, Thomas of an order that cannot be
escaped, and Bl.\ Ildefonso Schuster, commenting on this Mass, of malice that
Providence turns to its own ends; the three formulas agree on the one point the
reading needs. In this sense the adversary is not a second power. He is a
creature whose hatred is real and whose reach is lent.

\subsection{The enemy behind the enemy: St John Chrysostom}

If the war is real and bounded, who is the enemy? The Epistle answers, and St
John Chrysostom presses the answer hard. Paul speaks of ``wiles'' because the
devil ``never proposes to us sins in their proper colors''; he fights ``not
simply, nor directly, but with much wiliness''. The rulers of this darkness are
incorporeal, and the Christian overcomes them only by becoming what they are not:
``How then, you may say, are we to wrestle with the darkness? By becoming light.
How with the `spiritual hosts of wickedness'? By becoming good.'' Then Chrysostom
gives the sentence its sharpest point. Paul, he says, wrote to disciples who
would have many persecutors, and told them in effect, ``imagine not \dots\ that
it is they who war with you. They that really war with you, are the spirits that
work in them. Against them is our conflict'' (\work{Homilies on Ephesians} 22).
The persecutor is the instrument, and the enemy is the spirit that moves him.

From this Chrysostom draws a rule that he states with deliberate violence: ``If
any one hath an enemy, if any one hath been wronged by him, if any one is
exasperated, let him collect together all that wrath, all that fierceness, and
pour it out upon the head of the devil. Here wrath is a good thing, here anger is
profitable, here revenge is praiseworthy.'' ``Hath any one struck thee? Bear
malice against the devil, and never relinquish thy hatred towards him'' (Hom.\
22). The wrath is not abolished but redirected, from the brother who wronged
you to the adversary who moved him. The same Father, preaching on the parable of
this Mass's Gospel, makes the turn again and joins it to the Offertory's Job: ``For
what can be more wicked than the devil; yet nevertheless, even hence have we a
great opportunity of approving ourselves; and Job showeth it. But if the devil
hath become a cause of crowns, why art thou afraid of a man as an enemy?''
(\work{Homilies on Matthew} 61.5). The devil, who asked for Job's ruin, became the
occasion of his crown. A man who wrongs us can do no worse.

On the gospel of peace with which the feet are shod Chrysostom adds the paradox
that holds the reading together: ``this war which we have against them, puts an
end to another war, that, namely, which is between us and God; if we are at war
with the devil, we are at peace with God'' (\work{Homilies on Ephesians} 24). The
Christian who stands against the adversary stands with God; the one who makes his
brother the enemy has mistaken the field.

St Thomas, answering the objection that men plainly do persecute us, which he
draws from another verse of the same psalm, \latin{Multi qui persequuntur me}
(Ps 118:157), allows that they do, but says that flesh and blood attack us not
principally but as moved by a higher agent, \latin{a superiore movente, sc.\ a
diabolo} (\work{Super Epistolam ad Ephesios} 6, lect.~3). Chrysostom and Thomas
agree that the spirits are the real enemy; Thomas adds that men are their
instruments, which keeps the persecutor answerable for what he does.

\subsection{The armour that is Christ}

The armour is \latin{armatura Dei}, God's armour, and the lesson opens with
\latin{Confortamini in Domino, et in potentia virtutis eius}: the strength is
lent as the adversary's power is lent. St Jerome draws the conclusion that the
armour is the Lord himself. To put on the whole armour of God, he says, is to put
on Christ: \latin{quasi dixerit, Induite Dominum Jesum Christum}, ``But put ye on
the Lord Jesus Christ'' (Rom 13:14), so that the girdle of truth, the breastplate
of justice, the shield of faith and the rest are each Christ (\work{In Ephesios}
III, PL~26, col.~543). And he reads the wars of Israel against the nations,
David's against Goliath, Josue's and the Judges' and the Kings', as written for
us, a figure of the spiritual war. St Thomas names the pieces as the gifts and
virtues, the girdle temperance, the breastplate justice and the shield faith; the
continuation of \work{The Liturgical Year}, commenting on this Mass, holds that
the prophets describe the same pieces as Christ's own before they are ours:
``well, yes, but they are his, first'' (XI, p.~460). Jerome's identification
and Thomas's are different emphases, the one Christological and the other
moral. The soldier is clothed in Christ, and Christ in him becomes virtue.

The Gradual gives the soldier his fortress. St Robert Bellarmine explains that
the Hebrew for refuge ``signifies a well fortified house, placed on an
eminence'', where the tenants are secure from enemies, beasts, flood and storm,
and that those who dwell in God dwell ``in a city fortified by faith, hope, and
charity'' (\work{Explanation of the Psalms}, Ps~89:1). Durandus, expounding the
Gradual in the Mass of his Church, says that because even the well armed cannot
perfectly resist unless God helps them, the Church flees to God her refuge
(\work{Rationale} VI.138). The refuge that is God before the mountains
were made is the same will that the Introit says none resists: what threatens
the soldier lies inside it, and so does he.

\subsection{The persecutors of the Communion: St Hilary of Poitiers}

The Communion asks, \latin{quando facies de persequentibus me iudicium? iniqui
persecuti sunt me}. Sung after a Gospel that commands forgiveness from the heart,
the plea seems to ask God to punish the singer's enemies. St Hilary of Poitiers,
commenting on the verse, identifies the persecutors in the Epistle's terms. The
prophet, he says, desires that judgment be done on those who persecute him,
\latin{a diabolo scilicet, ministrisque ejus. Sunt enim spiritales nequitiae in
caelestibus}: on the devil, that is, and his ministers, for there are spirits of
wickedness in the heavenly places (Eph 6:12). What he asks is that, being made
wholly holy and pleasing to God, he may crush the devil within himself,
\latin{conterat intra se diabolum, conterat adversantes sibi inimicasque
virtutes, et fidei suae animae subdat nequitias spiritales}; not yet trusting his
own nature, he asks \latin{ut per misericordiam Dei judicium in inimicos suos
sit}. This is the power the Lord gave his apostles, \latin{Hanc potestatem judicii
Dominus apostolis suis dedit}, when he gave them power ``to tread upon serpents
and scorpions, and upon all the power of the enemy'' (Lk 10:19)
(\work{Tractatus in Psalmum} 118, \latin{Caph}~5, PL~9, cols.~573--574). Hilary
himself sets the Communion's verse beside the Epistle's. The judgment asked for
is not a curse on men but the power to tread the adversary down.

Hilary knows the last judgment and says so: \latin{Judicium post consummationem
saeculi Propheta esse non ignorat}, as Ps~1:5 shows; and he knows that the devil
himself is to be judged then, \latin{Sed scit et diabolum in resurrectionis
saeculo judicandum}. Why then does he take the prophet to be asking now for a
judgment within this life? He draws it from the first half of the verse, which
the Communion does not sing, ``How many are the days of thy servant?'': the
prophet knows himself to be a man of few days, \latin{cum se paucorum nunc dierum
secundum corporalem vitam sciat esse}, and so asks for what he may have in them.
The exposition of the psalm printed under the name of St Athanasius, on the same
unsung question, says likewise that the days of men on earth are few, in which
he prays that Satan be crushed at once under his feet (PG~27, col.~495; Rom
16:20). Hilary's identification of the persecutors and of the power asked stands
on the words the Communion sings; his ``now'' rests on the half it leaves out.

\subsection{Every element in the war}

Each element of the formulary takes its place in this reading. The Introit is
the soldier's confession before the battle: Mardochai, under sentence and
refusing to give a man the honour due to God, confesses the will that holds
all things. St Rabanus Maurus reads the edict of the heavenly king in Esther as
ruling the whole world and overcoming the pride of its persecutors, and Durandus
hears in the Introit the Church in Mardochai's voice, freed from the hand of the
devil as Mardochai was from Aman's. Its verse, ``Blessed are the undefiled in the
way, who walk in the law of the Lord'', names the road on which the soldier keeps
his footing; St Augustine finds the way in the law of the Lord. The Collect asks
the protection that Gregory's bridle supplies: that the family be free \latin{a
cunctis adversitatibus, te protegente}, and Schuster reads it as a prayer against
the obstacles the devil raises to turn men from good works. The Epistle names
the enemy and the armour, and the Gradual the fortress.

The Alleluia sings the first deliverance from an oppressor, Israel brought out
from a barbarous people. Cassiodorus reads the going out as release from the
host of the Egyptians, \latin{id est a daemonum turba}, and St John Chrysostom
on the same verse observes that God first freed his people from a cruel enemy and
only then asked their worship. The Gospel, which the next reading takes as its
centre, enters here at its end: the servant who will not forgive is handed
\latin{tortoribus}, and St Thomas glosses the torturers as the demons to whom he
is subjected, \latin{subiicitur daemonibus} (\work{Super Matthaeum} 18). The
fellow servant who owed a hundred pence was never the enemy; Chrysostom's
question, why fear a man as an enemy when the devil himself has become a cause of
crowns, is asked of this parable.

The Offertory is the reading's centre: the petition, the permission, the
measured loss. Rupert of Deutz, commenting on the Offertory in the Mass he
expounds, says that what happened to Job in history happens in the moral life to
every soldier of the virtues, whom the tempter strikes when a hidden
dispensation of God permits him (\work{De divinis officiis} XII.21); Durandus, that the devil is conquered above
all by patience (VI.138). The Secret offers the sacrifice by which God willed
to be appeased, and Chrysostom's sentence on the gospel of peace explains why
that offering belongs in a war: to be at war with the devil is to be at peace
with God. The Communion, read with Hilary, asks for the power to crush the
adversary. The Postcommunion names what the soldier has received, the food of
immortality, and asks that the pure mind follow what the mouth has taken; the
armour that is Christ is put on, in the end, by eating.

\subsection{The strongest difficulty}

The reading meets three objections. The first is that it makes God the author of
Job's ruin. Gregory and Thomas answer it by keeping the two wills apart: the
will to harm is the devil's own and evil, the permission is God's and good, and
God permits what the devil asks only so far as it serves the proving of his
servant. Gregory's \latin{ut deficeret} and \latin{ut proficeret} carry the whole
answer: the same trial is the devil's attempt at Job's fall and God's means of his
advance.

The second is that Chrysostom's redirection excuses the persecutor. If the
spirits are the real enemy, is the man who strikes me innocent? Chrysostom does
not say so. He tells the wronged man where to direct his wrath, not that the
wrong was none, and in the same homily on the parable he says that the one who
injures us harms himself. St Thomas keeps the persecutor's part by calling men
instruments moved by the devil, not principal agents but agents still. And
Chrysostom's ``here revenge is praiseworthy'' is aimed at the devil alone; it
licenses no vengeance on any man, and the next reading shows how hard the Gospel
presses that limit.

The third concerns the judgment itself. Hilary asks for power in this life to
crush the devil and his ministers, and says that what the prophet now asks is not
the judgment he knows is to come at the resurrection, \latin{non illud utique quod
in resurrectionis saeculo meminit futurum}. St Augustine hears in the same verse
the martyrs' cry for that final judgment, when vengeance falls on the Church's
persecutors; St Robert Bellarmine hears a plea for the condemnation of the wicked
who persecute and tempt. The Fathers differ in what is asked, a present power or
the last judgment, and in whom the plea names, not only in the time. Of the
others only Cassiodorus names Hilary's persecutors, the devil with his ministers,
and he joins them to the judgment at the end of the world. The answers are
compatible, not identical: Hilary affirms the judgment at the end, and the men
whom Augustine and Bellarmine see persecuting are, in Thomas's terms above,
instruments moved by the devil. The comparison at the end of the study sets out
where each finds his time.

\subsection{The four senses of this reading}

\begin{fourSenses}
\item[Literal.] Mardochai, under Aman's sentence, confesses that all things are
placed in God's will and that no one can resist it. Paul, a prisoner, arms the
Ephesians against principalities and powers and spirits of wickedness, not
against flesh and blood. The Book of Job tells that Satan asked leave to try a
just man and was given power over his goods and his flesh, and no further. The
psalmist, hoping in God's word, asks when God will judge his persecutors and begs
his help.
\item[Allegorical.] The armour of God is Christ, whom the Christian puts on
(Jerome); the wars of Israel are figures of the Church's spiritual war; Job,
whose name Gregory reads as ``grieving'', bears the type of the Redeemer together
with his Body (\work{Moralia}, Preface 14 and 16); the people under sentence
whom God delivers from Aman is the Church delivered from the devil (Rabanus,
Durandus).
\item[Moral.] Stand, and do not mistake the brother or even the persecutor for
the enemy; pour the wrath on the devil, and on him alone (Chrysostom). Take every
trial as measured by the hand that bridles the tempter, and turn it by patience
from the devil's attempt at your fall into God's means of your advance (Gregory,
Durandus).
\item[Anagogical.] The adversary, bridled now, is to be judged at the
resurrection, as Hilary affirms; and the soldier who has stood in the evil day
is crowned, the devil having become, in Chrysostom's word, a cause of crowns.
\end{fourSenses}
